\documentclass[10pt,a4paper]{amsart}
\usepackage[margin=19mm]{geometry}
\usepackage{amsmath,amssymb,mathtools}
\usepackage[scr=rsfs,cal=euler]{mathalfa}
\usepackage{xfrac}
\usepackage[unicode,hidelinks,bookmarks=false]{hyperref}
\newcommand{\ie}{i.e.\ }
\newcommand{\cf}{cf.\ }
\newcommand{\resp}{respectively\ }
\newcommand{\Subsection}{\subsection}
\providecommand{\nobreakdash}{\nobreak\hbox{-}}
\setlength{\emergencystretch}{2em}
\multlinegap0pt
\usepackage{bm}
\usepackage{bbm}
\usepackage{dsfont}
\usepackage{xspace}
\usepackage{cancel}
\usepackage{mathtools}
\usepackage{amsthm}
\makeatletter
\@ifundefined{lemma}{\newtheorem{lemma}{Lemma}}{}
\@ifundefined{proposition}{\newtheorem{proposition}[lemma]{Proposition}}{}
\makeatother
\newcommand{\Xset}{\mathcal{X}}
\title[On maximal curves of $n$-correct sets]{On maximal curves of $n$-correct sets}
\author{H. Hakopian, G. Vardanyan, N. Vardanyan}
\date{Source: arXiv:2507.11207v1 (CC BY 4.0)}
\begin{document}
\maketitle

\noindent\emph{Source and licence.} On maximal curves of $n$-correct sets. Version arXiv:2507.11207v1 (CC BY 4.0), distributed under the Creative Commons Attribution 4.0 International licence, as recorded in the arXiv licence metadata for this exact version and shown on the abstract page https://arxiv.org/abs/2507.11207v1. Full rights notice: licenses/arxiv-2507.11207-CC-BY-4.0.txt.\par\smallskip

\noindent\textit{Editorial scope.} Counted unit: the Proposition whose source label is \texttt{prop1011} in the article (source0.tex lines 614--631, source label \texttt{prop1011}), delivered together with the article's own complete original proof of it (source0.tex lines 632--680, one \texttt{proof} environment of 49 lines). No mathematics is written, repaired or re-derived: statement and proof are the article's own text line for line. Required context retained uncounted: dnk2 (lines 266--266); zsss2 (lines 451--452). Every definition, display and label of those lines is retained unchanged. Omitted from this package and not redistributed: 1--265, 267--450, 453--613, 681--890. Nothing inside a retained excerpt is omitted or altered: every retained body is delivered exactly as the article's lines are written, so no omission receipt of any kind accompanies this package. Citations: the retained excerpts carry no citation command at all, so no bibliography entry of the article is transcribed and no reference is redistributed. Reference adaptation: none was needed and none was made; every reference inside the delivered unit resolves inside it, so no unresolved reference or citation is produced. Printed slips kept literal: no printed word, symbol, hypothesis or number of the counted statement or of its proof was altered. Layout and class: the article's own class is \texttt{article} with packages including amssymb, amsthm, amsmath, graphicx, psfrag, array, hhline, enumitem; the wrapper is the lane's pinned \texttt{amsart} editorial wrapper at 10pt on A4, which restates the theorem-like environments the retained text needs and adds the article's own macro definitions that the retained text needs (\texttt{\textbackslash Xset}), with extra wrapper packages (bm, bbm, dsfont, xspace, cancel, mathtools). The delivered numbering is the unit's own. Assets: no retained span contains a figure environment, a graphics-inclusion command, a PDF-page inclusion, a table or a TikZ picture; no image file is copied into this package, the package declares no asset dependency and no image file is redistributed with it. Licence: version arXiv:2507.11207v1 of this article is distributed under the Creative Commons Attribution 4.0 International licence, as recorded in the arXiv licence metadata for this exact version and shown on the abstract page https://arxiv.org/abs/2507.11207v1. A full rights notice is delivered at licenses/arxiv-2507.11207-CC-BY-4.0.txt. Characters: every retained excerpt is pure ASCII, so the delivered engine, XeLaTeX with its default Latin Modern text font, reports no missing character.
\begin{equation}\label{dnk2}d(n, k) = (n-k+2)+(n-k+3)+\cdots+(n+1)\ \left[={k(2n+3-k)\over 2}\right].\end{equation}

\begin{equation}\label{zsss2} H_{k,m}^n = N_n-N_{n-k}-N_{n-m}+N_{n-k-m}. 
\end{equation}

\begin{proposition} \label{prop1011}
Suppose that $\Xset$ is an $n$-correct set and $f_1, f_2$ and $f_3$ are
three maximal curves, such that every two of them have no common components. Suppose also that $deg f_i=k_i, \ i=1,2,3,$ and $\sigma\ge -1.$
Then we have the following five expressions for $\#(f_1\cap f_2\cap f_3 \cap \Xset):$
\begin{enumerate}
\item
$N-d(n,k_1)-d(n,k_2)-d(n,k_3)+H_{k_1,k_2}^n+H_{k_2,k_3}^n+H_{k_1,k_3}^n,$
\item
$N-N_{n-k_1}-N_{n-k_2}-N_{n-k_3}+N_{n-k_1-k_2}+N_{n-k_2-k_3}+N_{n-k_1-k_3},$
\item
$N-d(n-k_1, k_2)-d(n-k_2,k_3)-d(n-k_3,k_1),$
\item
${1\over 2} \sigma(\sigma+1)-N_{k_1+k_2-n-3}-N_{k_2+k_3-n-3}-N_{k_3+k_1-n-3},$
\item
${1\over 2} \sigma(\sigma+1),\ \hbox{provided that}\  k_i+k_j\le n+2\ \forall\ 1\le i<j\le 3.$ 
\end{enumerate}

\end{proposition}

\begin{proof}
First let us verify that 
\begin{equation}\label{n+1+}
	\Xset\subset f_1\cup f_2 \cup f_3. 
\end{equation} Indeed, assume conversely that $A\in \Xset\setminus f_1f_2f_3.$
Then $f_1f_2f_3$ divides $P^\star_{A,\Xset},$ which, in view of the condition $\sigma\ge -1,$ i.e., $k_1+k_2+k_3\ge n+1,$ is contradiction.

Now denote $\gamma=\#(f_1\cap f_2\cap f_3 \cap \Xset).$ 
In view of \eqref{n+1+} we have that
$$ N=\#[(f_1\cup f_2\cup f_3 \cup)\cap \Xset]=   \#(f_1 \cap \Xset)+ \#(f_2\cap \Xset)+ \#(f_3 \cap \Xset)$$ $$-
\#(f_1\cap f_2 \cap \Xset) - \#(f_2\cap f_3 \cap \Xset) - \#(f_1\cap f_3 \cap \Xset) +\gamma$$
$$=d(n,k_1)+d(n,k_2)+d(n,k_3)-H_{k_1,k_2}^n-H_{k_2,k_3}^n-H_{k_1,k_3}^n+\gamma.$$ This proves the item (i).

Then we continue by using the equality \eqref{zsss2} here:
$$N= d(n,k_1)+d(n,k_2)+d(n,k_3)-3N+2N_{n-k_1}+2N_{n-k_2}+2N_{n-k_3}$$ $$-N_{n-k_1-k_2}-N_{n-k_2-k_3}-N_{n-k_1-k_3}+\gamma.$$
Next, in view of the definition of $d(n,k),$ we have
$$N= (N-N_{n-k_1})+(N-N_{n-k_2})+(N-N_{n-k_3})-3N+2N_{n-k_1}+2N_{n-k_2}+2N_{n-k_3}$$
$$-N_{n-k_1-k_2}-N_{n-k_2-k_3}-N_{n-k_1-k_3}+\gamma$$ 
$$=N_{n-k_1}-N_{n-k_1-k_2}+N_{n-k_2}-N_{n-k_2-k_3}+N_{n-k_3}-N_{n-k_1-k_3}+\gamma$$
$$=d(n-k_1, k_2)+d(n-k_2,k_3)+d(n-k_3,k_1)+\gamma.$$
In the last equality, which proves the item (iii) we again used the definition of $d(n,k).$ Note that the previous equality proves the item (ii).

Now let us prove (v).
In view of the item (iii) and \eqref{dnk2} we have that
$$\#(f_1\cap f_2\cap f_3 \cap \Xset)$$ 
$$= N- {1\over 2} k_2(2n-2k_1-k_2+3)- {1\over 2} k_3(2n-2k_2-k_3+3)- {1\over 2} k_1(2n-2k_3-k_1+3)$$
$$= N-{1\over 2} [(k_1+k_2+k_3)(2n+3)]-{1\over 2}[k_1^2+k_2^2+k_3^2+2k_1k_2+2k_2k_3+2k_3k_1]$$ 
$$= N-{1\over 2} [(k_1+k_2+k_3)(2n+3)]-{1\over 2}[k_1+k_2+k_3]^2$$ 
$$= N-{1\over 2} [(k_1+k_2+k_3)(2n+3-k_1-k_2-k_3)]$$ 
$$= {1\over 2} [(n+2)(n+1)]-{1\over 2} [(n+2+\sigma)(n+1-\sigma)]= {1\over 2} \sigma(\sigma+1).$$ 
At the end let us prove the item (iv) which is an extension of the item (v).
For this end let us define the following quantity $\widetilde d(n,k)$ for all $n,k\ge 0:$ 
\begin{equation}\label{dnk3+}\widetilde d(n, k) := (n-k+2)+(n-k+3)+\cdots+(n+1)\ \left[={k(2n+3-k)\over 2}\right].\end{equation}

Let us verify that 
\begin{equation}\label{dnk3++}\widetilde d(n, k)=d(n,k)-N_{-n+k-3},\ \forall n,k\ge 0.\end{equation}
Indeed, if $0\le k\le n+2$ then, in view of \eqref{dnk2}, we have that $ \widetilde d(n, k)=d(n,k)=d(n,k)-N_{-n+k-3}.$
While if  $k\ge n+3$ then we have that $ \widetilde d(n, k)=N_n-N_{-n+k-3}=d(n,k)-N_{-n+k-3}.$

Next observe that the proof of the item (v), in view of \eqref{dnk3+}, means that
\begin{equation}\label{dnk3-}N-\widetilde d(n-k_1, k_2)-\widetilde d(n-k_2,k_3)-\widetilde d(n-k_3,k_1)={1\over 2} \sigma(\sigma+1),\end{equation}
where $0\le k_i\le n,\ i=1,2,3.$

\noindent Finally, by using the item (iii) and the 
relations \eqref{dnk3++} and \eqref{dnk3-} , we obtain that 
$$\#(f_1\cap f_2\cap f_3 \cap \Xset)=N-d(n-k_1, k_2)- d(n-k_2,k_3)-d(n-k_3,k_1)$$ $$=N-\widetilde d(n-k_1, k_2)-\widetilde d(n-k_2,k_3)-\widetilde d(n-k_3,k_1)$$ 
$$-N_{k_1+k_2-n-3}-N_{k_2+k_3-n-3}-N_{k_3+k_1-n-3}$$
$$={1\over 2} \sigma(\sigma+1)-N_{k_1+k_2-n-3}-N_{k_2+k_3-n-3}-N_{k_3+k_1-n-3}.$$
\end{proof}

\end{document}
